%%%%%%%%%%%%%%%%%%%%%%%%%%%%%%%%%%%%%%%%%%%%%%%%%%%%%%%%%%%%%%%%%%%%%%%%%%%%%%%
%% TKS FORMAL MATHEMATICAL MANUAL v6.0 — ACADEMIC RELEASE
%% Part 3 of 4: Category Theory, Denotational Semantics, Type Theory
%%%%%%%%%%%%%%%%%%%%%%%%%%%%%%%%%%%%%%%%%%%%%%%%%%%%%%%%%%%%%%%%%%%%%%%%%%%%%%%

%% ============================================================================
%% PART III: CATEGORICAL AND SEMANTIC FRAMEWORK
%% ============================================================================
\part{Categorical and Semantic Framework}

%% ============================================================================
%% CHAPTER 6: THE CATEGORY NOETICA
%% ============================================================================
\chapter{The Category Noetica}
\label{ch:category-noetica}

This chapter provides a rigorous category-theoretic foundation for TKS, proving that the system forms a valid category and exploring its subcategories.

\section{Definition of Category Noetica}

\begin{definition}[Category Noetica]
\label{def:category-noetica}
The category $\Noetica$ is defined by:

\textbf{Objects:}
\[
\Ob(\Noetica) = \Ideas \cup \Worlds \cup \Elements \cup \Foundations \cup \Acquisitions
\]

\textbf{Morphisms:}
\[
\Hom(\Noetica) = \Noetics \cup \{\oplus_W, \ominus_W : W \in \Worlds\}
\]

For objects $X, Y \in \Ob(\Noetica)$:
\[
\Hom(X, Y) = \{f : X \to Y \mid f \text{ is a valid TKS transformation}\}
\]

\textbf{Identity:}
\[
\id_X = \noe{0} : X \to X
\]

\textbf{Composition:}
For morphisms $f : X \to Y$ and $g : Y \to Z$:
\[
(g \circ f) : X \to Z, \quad (g \circ f)(x) = g(f(x))
\]
\end{definition}

\section{Verification of Category Axioms}

\begin{theorem}[Noetica is a Category]
\label{thm:noetica-category}
The structure $\Noetica$ satisfies all category axioms.
\end{theorem}

\begin{proof}
We verify each axiom:

\textbf{(C1) Identity Law:}
For any morphism $f : X \to Y$:
\begin{align}
f \circ \id_X &= f \circ \noe{0} = f \\
\id_Y \circ f &= \noe{0} \circ f = f
\end{align}
Both equalities hold because $\noe{0}(x) = x$ for all $x$ (Axiom~\ref{ax:identity}).

\textbf{(C2) Associativity:}
For morphisms $f : X \to Y$, $g : Y \to Z$, $h : Z \to W$:
\begin{align}
(h \circ (g \circ f))(x) &= h((g \circ f)(x)) = h(g(f(x))) \\
((h \circ g) \circ f)(x) &= (h \circ g)(f(x)) = h(g(f(x)))
\end{align}
Thus $h \circ (g \circ f) = (h \circ g) \circ f$ by Axiom~\ref{ax:assoc-comp}.

\textbf{(C3) Composition Closure:}
For any $f \in \Hom(X, Y)$ and $g \in \Hom(Y, Z)$, the composite $g \circ f$ is in $\Hom(X, Z)$ by closure of Noetic composition (Theorem~\ref{thm:closure}).
\end{proof}

\section{World Subcategories}

\begin{definition}[World Subcategory]
\label{def:world-subcat}
For each World $W \in \{A, B, C, D\}$, define the subcategory $\Noetica_W$:
\begin{align}
\Ob(\Noetica_W) &= \{X \in \Ob(\Noetica) : \mathrm{World}(X) = W\} \\
\Hom(\Noetica_W) &= \{f \in \Hom(\Noetica) : f \text{ preserves World } W\}
\end{align}
\end{definition}

\begin{proposition}[World Subcategories are Categories]
\label{prop:world-subcat}
Each $\Noetica_W$ is a valid subcategory of $\Noetica$.
\end{proposition}

\begin{proof}
\begin{enumerate}
  \item Identity: $\id_X = \noe{0}$ preserves World($X$) = $W$.
  \item Composition: If $f, g$ both preserve World $W$, so does $g \circ f$.
  \item Closure: The restriction to World-$W$ objects and morphisms is closed.
\end{enumerate}
\end{proof}

\begin{definition}[Element Subcategory]
\label{def:element-subcat}
The \textbf{Element subcategory} $\Noetica_\Elements$ is:
\begin{align}
\Ob(\Noetica_\Elements) &= \Elements \quad \text{(40 Elements)} \\
\Hom(\Noetica_\Elements) &= \Noetics \quad \text{(10 Noetic operators)}
\end{align}
\end{definition}

%% ============================================================================
%% CHAPTER 7: THE ACBE FUNCTOR
%% ============================================================================
\chapter{The ACBE Functor}
\label{ch:acbe-functor}

The ACBE (Above-Cause-Below-Effect) transformation is the primary manifestation functor in TKS, mapping spiritual causes to physical effects.

\section{Functor Definition}

\begin{definition}[ACBE Functor]
\label{def:acbe-functor}
The \textbf{ACBE functor} is:
\[
\mathrm{ACBE} : \Noetica_A \to \Noetica_D
\]

\textbf{Object Mapping:}
For each $An \in \Ob(\Noetica_A)$:
\[
\mathrm{ACBE}(An) = Dn
\]

Explicitly:
\begin{center}
\begin{tabular}{ll}
$\mathrm{ACBE}(A1) = D1$ & Divine Mind $\to$ Physical Brain \\
$\mathrm{ACBE}(A2) = D2$ & Spiritual Positive $\to$ Physical Order \\
$\mathrm{ACBE}(A3) = D3$ & Spiritual Negative $\to$ Physical Disorder \\
$\mathrm{ACBE}(A4) = D4$ & Aetheric Vibration $\to$ Physical Vibration \\
$\mathrm{ACBE}(A5) = D5$ & Soul-Womb $\to$ Physical Vessel \\
$\mathrm{ACBE}(A6) = D6$ & Spiritual Discipline $\to$ Physical Delivery \\
$\mathrm{ACBE}(A7) = D7$ & Destiny Pattern $\to$ Physical Rhythm \\
$\mathrm{ACBE}(A8) = D8$ & Esoteric Truth $\to$ High Quality \\
$\mathrm{ACBE}(A9) = D9$ & Exoteric Symbol $\to$ Low Quality \\
$\mathrm{ACBE}(A10) = D10$ & Akashic Pattern $\to$ Physical Form \\
\end{tabular}
\end{center}

\textbf{Morphism Mapping:}
For morphism $f : X \to Y$ in $\Noetica_A$:
\[
\mathrm{ACBE}(f) : \mathrm{ACBE}(X) \to \mathrm{ACBE}(Y)
\]
\[
\mathrm{ACBE}(\noe{k}) = \noe{k}
\]
Noetics are preserved across worlds.
\end{definition}

\section{Proof of Functoriality}

\begin{theorem}[ACBE is a Functor]
\label{thm:acbe-functor}
ACBE preserves identity and composition, hence is a valid functor.
\end{theorem}

\begin{proof}
\textbf{(F1) Identity Preservation:}
\[
\mathrm{ACBE}(\id_X) = \mathrm{ACBE}(\noe{0}) = \noe{0} = \id_{\mathrm{ACBE}(X)}
\]

\textbf{(F2) Composition Preservation:}
For $f : X \to Y$ and $g : Y \to Z$ in $\Noetica_A$:
\begin{align}
\mathrm{ACBE}(g \circ f) &= \mathrm{ACBE}(g) \circ \mathrm{ACBE}(f)
\end{align}
Since $\mathrm{ACBE}(\noe{j} \circ \noe{i}) = \noe{j} \circ \noe{i} = \mathrm{ACBE}(\noe{j}) \circ \mathrm{ACBE}(\noe{i})$.
\end{proof}

\section{The Cascade Decomposition}

\begin{theorem}[ACBE Factorization]
\label{thm:acbe-factor}
ACBE factors through intermediate categories:
\[
\mathrm{ACBE} = F_D \circ F_C \circ F_B
\]
where:
\begin{align}
F_B &: \Noetica_A \to \Noetica_B && \text{(Spiritual } \to \text{ Mental)} \\
F_C &: \Noetica_B \to \Noetica_C && \text{(Mental } \to \text{ Emotional)} \\
F_D &: \Noetica_C \to \Noetica_D && \text{(Emotional } \to \text{ Physical)}
\end{align}
\end{theorem}

\begin{proof}
Each $F_W$ is defined by:
\begin{itemize}
  \item $F_B(An) = Bn$, $F_B(\noe{k}) = \noe{k}$
  \item $F_C(Bn) = Cn$, $F_C(\noe{k}) = \noe{k}$
  \item $F_D(Cn) = Dn$, $F_D(\noe{k}) = \noe{k}$
\end{itemize}
Then $(F_D \circ F_C \circ F_B)(An) = F_D(F_C(Bn)) = F_D(Cn) = Dn = \mathrm{ACBE}(An)$.
\end{proof}

\begin{remark}[Commutative Diagram]
The ACBE cascade forms a commutative diagram:
\[
\begin{tikzcd}
\Noetica_A \arrow[r, "F_B"] \arrow[rrr, bend left=20, "\mathrm{ACBE}"] & \Noetica_B \arrow[r, "F_C"] & \Noetica_C \arrow[r, "F_D"] & \Noetica_D
\end{tikzcd}
\]
\end{remark}

%% ============================================================================
%% CHAPTER 8: RPM AS A MONAD
%% ============================================================================
\chapter{RPM as a Recursive Prerequisite Monad}
\label{ch:rpm-monad}

The Recursive Prerequisite Model (RPM) is formalized as a monad on the category of Acquisitions.

\section{The Acquisition Category}

\begin{definition}[Prerequisite Operators]
\label{def:prereq-ops}
The \textbf{Prerequisite operator set} is:
\[
\Pi = \{D, W, P\}
\]
where:
\begin{itemize}
  \item $D$: Desire --- initiates and directs
  \item $W$: Wisdom --- models and validates
  \item $P$: Power --- executes and manifests
\end{itemize}
\end{definition}

\begin{theorem}[Strict Ordering]
\label{thm:prereq-order}
\[
D \prec W \prec P
\]
Desire must precede Wisdom must precede Power.
\end{theorem}

\begin{definition}[Pure Desire]
\label{def:pure-desire}
\textbf{A0} (Pure Desire) is the root intentional vector:
\[
A0 = \mathrm{PureDesire}
\]
the pre-foundational desire before attachment to any specific Foundation.
\end{definition}

\begin{definition}[Acquisition Space]
\label{def:acq-space}
The \textbf{Acquisition space} is:
\[
\Acquisitions = \{A0\} \cup \{D_n : n \in \{1,\ldots,7\}\} \cup \{W_n : n \in \{1,\ldots,7\}\} \cup \{P_n : n \in \{1,\ldots,7\}\}
\]
with $|\Acquisitions| = 1 + 7 + 7 + 7 = 22$.
\end{definition}

\begin{definition}[Category Acq]
\label{def:cat-acq}
The \textbf{category Acq} is:
\begin{align}
\Ob(\Acq) &= \Acquisitions \\
\Hom(\Acq) &= \{\to\} \quad \text{(dependency relation)}
\end{align}
For $X, Y \in \Acquisitions$:
\[
X \to Y \in \Hom(X, Y) \Leftrightarrow Y \text{ depends on } X
\]
\end{definition}

\section{The RPM Endofunctor}

\begin{definition}[RPM Endofunctor]
\label{def:rpm-endofunctor}
The \textbf{RPM endofunctor} $\RPM : \Acq \to \Acq$ is defined by:

\textbf{Object Mapping:}
\[
\RPM(X) = \begin{cases}
X & \text{if } \mathrm{Satisfied}(X) \\
\bot & \text{if } \neg\mathrm{Satisfied}(X)
\end{cases}
\]

\textbf{Morphism Mapping:}
\[
\RPM(X \to Y) = \begin{cases}
X \to Y & \text{if } \mathrm{Satisfied}(X) \\
\bot \to \bot & \text{otherwise}
\end{cases}
\]
\end{definition}

\begin{definition}[Satisfaction Predicate]
\label{def:satisfied}
The predicate $\mathrm{Satisfied} : \Acquisitions \to \{\mathtt{true}, \mathtt{false}\}$ determines whether an acquisition is currently met.
\end{definition}

\section{Monad Structure}

\begin{definition}[Unit Natural Transformation]
\label{def:rpm-unit}
The \textbf{unit} $\eta : \mathrm{Id}_{\Acq} \Rightarrow \RPM$ is:
\[
\eta_X : X \to \RPM(X)
\]
\[
\eta_X = \begin{cases}
\id_X & \text{if } \mathrm{Satisfied}(X) \\
\bot_X & \text{otherwise}
\end{cases}
\]
\end{definition}

\begin{definition}[Multiplication Natural Transformation]
\label{def:rpm-mult}
The \textbf{multiplication} $\mu : \RPM^2 \Rightarrow \RPM$ is:
\[
\mu_X : \RPM(\RPM(X)) \to \RPM(X)
\]
\[
\mu_X = \id_{\RPM(X)}
\]
(Flattening is trivial for this monad.)
\end{definition}

\begin{theorem}[RPM Satisfies Monad Laws]
\label{thm:rpm-monad-laws}
The structure $(\RPM, \eta, \mu)$ satisfies the three monad laws.
\end{theorem}

\begin{proof}
\textbf{(M1) Left Unit Law:} $\mu \circ \RPM\eta = \id_\RPM$

For any $X$:
\[
(\mu \circ \RPM\eta)_X = \mu_X \circ \RPM(\eta_X) = \id_{\RPM(X)} \circ \RPM(\eta_X) = \RPM(\eta_X)
\]
If $\mathrm{Satisfied}(X)$, then $\eta_X = \id_X$, so $\RPM(\eta_X) = \id_{\RPM(X)}$.
If not, both sides yield $\bot$.

\textbf{(M2) Right Unit Law:} $\mu \circ \eta\RPM = \id_\RPM$

Similar reasoning.

\textbf{(M3) Associativity:} $\mu \circ \RPM\mu = \mu \circ \mu\RPM$

Both compositions flatten nested RPM applications identically.
\end{proof}

\section{Monadic Bind Operation}

\begin{definition}[Kleisli Bind]
\label{def:kleisli-bind}
The \textbf{bind operation} is:
\[
(\gg\!=) : \RPM(X) \times (X \to \RPM(Y)) \to \RPM(Y)
\]
\[
m \gg\!= f = \begin{cases}
f(x) & \text{if } m = \mathrm{Success}(x) \\
\bot & \text{if } m = \bot
\end{cases}
\]
\end{definition}

\begin{example}[Chain Evaluation]
\label{ex:chain-eval}
Evaluating the prerequisite chain for Foundation $n$:
\begin{align}
\texttt{evalChain}(n) &= \RPM(A0) \gg\!= \lambda a_0. \\
&\quad\; \RPM(D_n) \gg\!= \lambda d_n. \\
&\quad\; \RPM(W_n) \gg\!= \lambda w_n. \\
&\quad\; \RPM(P_n) \gg\!= \lambda p_n. \\
&\quad\; \mathrm{return}(\mathrm{Outcome}_n)
\end{align}
\end{example}

\section{RPM Diagnostic Algorithm}

\begin{definition}[RPM Diagnostic Function]
\label{def:rpm-diag}
\begin{verbatim}
RPM : Acquisitions × Goal → (Status, FailureOrigin?)

RPM(X, G) =
  if X = G and Satisfied(X):
    return (SUCCESS, null)
  elif not Satisfied(X):
    return (FAILURE, X)
  else:
    for each Y ∈ Dependencies(X):
      (status, origin) = RPM(Y, G)
      if status = FAILURE:
        return (FAILURE, origin)
    return (SUCCESS, null)
\end{verbatim}
\end{definition}

\begin{theorem}[RPM Termination]
\label{thm:rpm-term}
The RPM diagnostic algorithm terminates for any input.
\end{theorem}

\begin{proof}
\textbf{(1) Finiteness:} $|\Acquisitions| = 22$.

\textbf{(2) Acyclicity:} The dependency graph is a DAG by Theorem~\ref{thm:prereq-order}. For any Foundation $n$:
\[
A0 \to D_n \to W_n \to P_n \to \mathrm{Outcome}_n
\]
No backward edges exist because $D \prec W \prec P$.

\textbf{(3) Termination:} Each node is visited at most once. Maximum path length is 4. Total complexity is $O(|\Acquisitions|) = O(22)$.
\end{proof}

%% ============================================================================
%% CHAPTER 9: DENOTATIONAL SEMANTICS
%% ============================================================================
\chapter{Formal Semantics of TKS}
\label{ch:denotational-semantics}

This chapter provides complete denotational semantics for TKS, giving formal mathematical meaning to all operations.

\section{Syntax vs. Semantics}

\begin{definition}[Syntactic Domain]
\label{def:syntactic-domain}
The \textbf{syntactic domain} $\mathrm{Syn}$ consists of:
\begin{itemize}
  \item Expression syntax (defined by BNF grammar)
  \item Operator symbols ($\noe{k}$, $\Tadd$, etc.)
  \item Structural elements (Elements, Worlds, Foundations)
\end{itemize}
\end{definition}

\begin{definition}[Semantic Domain]
\label{def:semantic-domain}
The \textbf{semantic domain} $\mathcal{D}$ consists of:
\begin{itemize}
  \item $\mathcal{D}_\Ideas$: The space of Idea values (graded multisets)
  \item $\mathcal{D}_\mathbb{B}$: Boolean values $\{\mathtt{true}, \mathtt{false}\}$
  \item $\mathcal{D}_\bot$: Error values $\{\bot\}$
  \item $\mathcal{D} = \mathcal{D}_\Ideas \cup \mathcal{D}_\mathbb{B} \cup \mathcal{D}_\bot$
\end{itemize}
\end{definition}

\begin{definition}[Semantic Function]
\label{def:sem-function}
The \textbf{semantic function} $\sem{\cdot}$ maps syntax to semantics:
\[
\sem{\cdot} : \mathrm{Syn} \to \mathcal{D}
\]
\end{definition}

\section{Semantics of Base Expressions}

\begin{definition}[Element Semantics]
\label{def:elem-sem}
For Element $Wn$ where $W \in \Worlds$ and $n \in \{1,\ldots,10\}$:
\[
\sem{Wn} = e_{W,n} \in \mathcal{D}_\Ideas
\]
where $e_{W,n}$ is the canonical Idea value for that Element.
\end{definition}

\begin{definition}[World Semantics]
\label{def:world-sem}
For World $W$:
\[
\sem{W} = \{e_{W,n} : n \in \{1,\ldots,10\}\} \subseteq \mathcal{D}_\Ideas
\]
\end{definition}

\section{Semantics of Noetic Operators}

\begin{definition}[Noetic Operator Semantics]
\label{def:noetic-sem}
For each Noetic $\noe{k}$:
\[
\sem{\noe{k}} : \mathcal{D}_\Ideas \to \mathcal{D}_\Ideas
\]
is a total function defined by:
\[
\sem{\noe{k}}(d) = \mathfrak{n}_k(d)
\]
where $\mathfrak{n}_k$ is the denotation of Noetic $k$.
\end{definition}

\begin{definition}[Concrete Noetic Denotations]
\label{def:noetic-denote}
The concrete denotations are:
\begin{align}
\mathfrak{n}_0(d) &= d && \text{(Identity)} \\
\mathfrak{n}_1(d) &= \mathrm{Aware}(d) && \text{(Apply awareness)} \\
\mathfrak{n}_2(d) &= \mathrm{Attract}(d) && \text{(Apply attraction)} \\
\mathfrak{n}_3(d) &= \mathrm{Repel}(d) && \text{(Apply repulsion)} \\
\mathfrak{n}_4(d) &= \mathrm{Vibrate}(d) && \text{(Apply vibration)} \\
\mathfrak{n}_5(d) &= \mathrm{Receive}(d) && \text{(Apply reception)} \\
\mathfrak{n}_6(d) &= \mathrm{Project}(d) && \text{(Apply projection)} \\
\mathfrak{n}_7(d) &= \mathrm{Rhythm}(d) && \text{(Apply periodicity)} \\
\mathfrak{n}_8(d) &= \mathrm{Elevate}(d) && \text{(Apply elevation)} \\
\mathfrak{n}_9(d) &= \mathrm{Ground}(d) && \text{(Apply grounding)}
\end{align}
\end{definition}

\begin{definition}[Noetic Application Semantics]
\label{def:noetic-app-sem}
For expression $E^k$:
\[
\sem{E^k} = \sem{\noe{k}}(\sem{E}) = \mathfrak{n}_k(\sem{E})
\]
\end{definition}

\section{Semantics of Tootra Arithmetic}

\begin{definition}[Tootra-Addition Semantics]
\label{def:tadd-sem}
\[
\sem{X \Tadd Y} = \sem{X} \uplus \sem{Y}
\]
where $\uplus$ is multiset union.
\end{definition}

\begin{definition}[Tootra-Subtraction Semantics]
\label{def:tsub-sem}
\[
\sem{X \Tsub Y} = \sem{X} \setminus_m \sem{Y}
\]
where $\setminus_m$ is multiset difference.
\end{definition}

\begin{definition}[Tootra-Multiplication Semantics]
\label{def:tmul-sem}
\[
\sem{X \Tmul Y} = \sem{X} \times_H \sem{Y}
\]
where $\times_H$ is the Hadamard (component-wise) product.
\end{definition}

\begin{definition}[Tootra-Division Semantics]
\label{def:tdiv-sem}
\[
\sem{X \Tdiv Y} = \sem{X} /_m \sem{Y}
\]
where $/_m$ is component-wise floor division.
\end{definition}

\section{Semantics of Composition}

\begin{definition}[Composition Semantics]
\label{def:comp-sem}
For composed operations $f \circ g$:
\[
\sem{f \circ g} = \sem{f} \circ \sem{g}
\]
where the right-hand $\circ$ is function composition in $\mathcal{D}$.
\end{definition}

\begin{theorem}[Compositional Semantics]
\label{thm:compositional}
TKS semantics is compositional: the meaning of a complex expression is determined entirely by the meanings of its parts.
\end{theorem}

\begin{proof}
By structural induction on expressions:
\begin{itemize}
  \item Base case: Atoms (Elements, Worlds) have direct denotations.
  \item Inductive case: Compound expressions are defined in terms of their subexpressions' denotations.
\end{itemize}
\end{proof}

\section{Semantics of Fractals}

\begin{definition}[Fractal Semantics]
\label{def:fractal-sem}
For fractal $\langle k_1 : k_2 : \cdots : k_n \rangle(E)$:
\[
\sem{\langle k_1 : k_2 : \cdots : k_n \rangle(E)} = \mathfrak{n}_{k_1}(\mathfrak{n}_{k_2}(\cdots \mathfrak{n}_{k_n}(\sem{E})\cdots))
\]
Application is right-to-left (innermost first).
\end{definition}

\begin{theorem}[Fractal Correctness]
\label{thm:fractal-correct}
Fractal semantics correctly implements nested Noetic application.
\end{theorem}

\begin{proof}
By definition:
\[
\sem{\langle X : Y : Z \rangle(E)} = \sem{\noe{X}}(\sem{\noe{Y}}(\sem{\noe{Z}}(\sem{E}))) = \sem{E^{ZYX}}
\]
which matches the intended right-to-left application.
\end{proof}

\section{Semantics of ACBE}

\begin{definition}[ACBE Semantics]
\label{def:acbe-sem}
\[
\sem{\mathrm{ACBE}(E)} = \mathfrak{f}_D(\mathfrak{f}_C(\mathfrak{f}_B(\sem{E})))
\]
where $\mathfrak{f}_W$ is the world-transition function for world $W$.
\end{definition}

\begin{definition}[World Transition Functions]
\label{def:world-trans}
\begin{align}
\mathfrak{f}_B(d) &= \mathrm{Translate}_B(d) && \text{Spiritual} \to \text{Mental} \\
\mathfrak{f}_C(d) &= \mathrm{Translate}_C(d) && \text{Mental} \to \text{Emotional} \\
\mathfrak{f}_D(d) &= \mathrm{Translate}_D(d) && \text{Emotional} \to \text{Physical}
\end{align}
\end{definition}

\section{Semantics of RPM}

\begin{definition}[RPM Semantics]
\label{def:rpm-sem}
\[
\sem{\mathrm{RPM}(X, G)} = \begin{cases}
\mathtt{success} & \text{if } \mathrm{ChainValid}(\sem{X}, \sem{G}) \\
\mathtt{failure}(o) & \text{otherwise, with failure origin } o
\end{cases}
\]
\end{definition}

%% ============================================================================
%% CHAPTER 10: TYPE THEORY OF TKS
%% ============================================================================
\chapter{Type Theory of TKS Expressions}
\label{ch:type-theory}

This chapter presents a formal type system for TKS expressions with inference rules and type safety properties.

\section{Base Types}

\begin{definition}[TKS Base Types]
\label{def:base-types}
The base types are:
\begin{align}
\mathtt{World} &: \text{Type} && \{A, B, C, D\} \\
\mathtt{Mode} &: \text{Type} && \{1, 2, \ldots, 10\} \\
\mathtt{Foundation} &: \text{Type} && \{F_1, \ldots, F_7\} \\
\mathtt{SubWorld} &: \text{Type} && \{a, b, c, d\}
\end{align}
\end{definition}

\section{Compound Types}

\begin{definition}[Compound Type Signatures]
\label{def:compound-types}
\begin{align}
\mathtt{Element} &: \mathtt{World} \to \mathtt{Mode} \to \mathtt{Type} \\
\mathtt{Noetic} &: \mathtt{Expr} \to \mathtt{Expr} \\
\mathtt{SubFound} &: \mathtt{Foundation} \to \mathtt{SubWorld} \to \mathtt{Type} \\
\mathtt{Acquisition} &: \mathtt{Foundation} \to \mathtt{Bool} \\
\mathtt{Expression} &: \mathtt{Element} \times \mathtt{Noetic}^* \times \mathtt{SubFound}^? \to \mathtt{Idea} \\
\mathtt{Fractal} &: \mathtt{List}(\mathtt{Noetic}) \to \mathtt{Expr} \to \mathtt{Expr}
\end{align}
\end{definition}

\section{Type Inference Rules}

\begin{definition}[Type Judgment]
\label{def:type-judgment}
A type judgment has the form:
\[
\Gamma \types e : \tau
\]
meaning ``in context $\Gamma$, expression $e$ has type $\tau$.''
\end{definition}

\begin{definition}[Element Formation Rule]
\label{def:elem-form}
\[
\frac{W : \mathtt{World} \quad n : \mathtt{Mode}}{\Gamma \types Wn : \mathtt{Element}} \quad [\textsc{Elem-Form}]
\]
\end{definition}

\begin{definition}[Noetic Application Rule]
\label{def:noetic-rule}
\[
\frac{\Gamma \types E : \mathtt{Expr} \quad k \in \{0,\ldots,9\}}{\Gamma \types E^k : \mathtt{Expr}} \quad [\textsc{Noetic-Apply}]
\]
\end{definition}

\begin{definition}[Foundation Restriction Rule]
\label{def:found-rule}
\[
\frac{\Gamma \types E : \mathtt{Expr} \quad m \in \{1,\ldots,7\} \quad w \in \{a,b,c,d\}}{\Gamma \types E_{mw} : \mathtt{Expr}} \quad [\textsc{Found-Restrict}]
\]
\end{definition}

\begin{definition}[World Association Rule]
\label{def:world-rule}
\[
\frac{\Gamma \types E : \mathtt{Expr} \quad W : \mathtt{World}}{\Gamma \types E \oplus W : \mathtt{Expr}} \quad [\textsc{World-Assoc}]
\]
\end{definition}

\begin{definition}[Tootra Addition Rule]
\label{def:tadd-rule}
\[
\frac{\Gamma \types X : \mathtt{Idea} \quad \Gamma \types Y : \mathtt{Idea}}{\Gamma \types X \Tadd Y : \mathtt{Idea}} \quad [\textsc{T-Add}]
\]
\end{definition}

\begin{definition}[Fractal Rule]
\label{def:fractal-rule}
\[
\frac{\Gamma \types E : \mathtt{Expr} \quad \forall i. k_i \in \{0,\ldots,9\} \quad n \geq 1}{\Gamma \types \langle k_1 : \cdots : k_n \rangle(E) : \mathtt{Expr}} \quad [\textsc{Fractal}]
\]
\end{definition}

\section{Type Compatibility}

\begin{definition}[Cross-World Compatibility]
\label{def:cross-world}
\[
\frac{\Gamma \types E : \mathtt{Expr}_{W_1} \quad W_2 : \mathtt{World} \quad d(W_1, W_2) \leq 3}{\Gamma \types E \oplus W_2 : \mathtt{Expr}_{W_2}} \quad [\textsc{Cross-World}]
\]
\end{definition}

\section{Error Types}

\begin{definition}[Type Error Categories]
\label{def:error-types}
\begin{center}
\begin{tabular}{ll}
\toprule
\textbf{Error Type} & \textbf{Description} \\
\midrule
$\mathtt{TypeError}$ & Malformed expression syntax \\
$\mathtt{DomainError}$ & Operation outside valid domain \\
$\mathtt{WorldError}$ & Invalid cross-world operation \\
$\mathtt{NoeticError}$ & Invalid Noetic application \\
$\mathtt{FoundationError}$ & Invalid Foundation context \\
$\mathtt{RPMError}$ & Prerequisite chain failure \\
$\mathtt{FractalError}$ & Invalid fractal structure \\
\bottomrule
\end{tabular}
\end{center}
\end{definition}

\begin{definition}[Illegal Combinations]
\label{def:illegal}
The following are ill-typed:
\begin{itemize}
  \item $\mathtt{Element}(W, 11)$ --- Mode out of range
  \item $A8_{8a}$ --- Foundation 8 doesn't exist
  \item $(A8 \oplus D)^{10}$ --- Noetic 10 doesn't exist
  \item $\langle\rangle(E)$ --- Empty fractal
\end{itemize}
\end{definition}

\section{Type Safety}

\begin{theorem}[Progress]
\label{thm:progress}
If $\Gamma \types e : \tau$ and $e$ is not a value, then there exists $e'$ such that $e \to e'$.
\end{theorem}

\begin{theorem}[Preservation]
\label{thm:preservation}
If $\Gamma \types e : \tau$ and $e \to e'$, then $\Gamma \types e' : \tau$.
\end{theorem}

\begin{corollary}[Type Safety]
\label{cor:type-safety}
Well-typed TKS expressions do not get stuck (they either evaluate to a value or continue reducing).
\end{corollary}

%% End of Part 3
%% Continue in Part 4 with Compiler Spec, Theorems, Examples, Appendices

\end{document}
